% Options for packages loaded elsewhere
\PassOptionsToPackage{unicode}{hyperref}
\PassOptionsToPackage{hyphens}{url}
%
\documentclass[
]{article}
\usepackage{amsmath,amssymb}
\usepackage{iftex}
\ifPDFTeX
  \usepackage[T1]{fontenc}
  \usepackage[utf8]{inputenc}
  \usepackage{textcomp} % provide euro and other symbols
\else % if luatex or xetex
  \usepackage{unicode-math} % this also loads fontspec
  \defaultfontfeatures{Scale=MatchLowercase}
  \defaultfontfeatures[\rmfamily]{Ligatures=TeX,Scale=1}
\fi
\usepackage{lmodern}
\ifPDFTeX\else
  % xetex/luatex font selection
\fi
% Use upquote if available, for straight quotes in verbatim environments
\IfFileExists{upquote.sty}{\usepackage{upquote}}{}
\IfFileExists{microtype.sty}{% use microtype if available
  \usepackage[]{microtype}
  \UseMicrotypeSet[protrusion]{basicmath} % disable protrusion for tt fonts
}{}
\makeatletter
\@ifundefined{KOMAClassName}{% if non-KOMA class
  \IfFileExists{parskip.sty}{%
    \usepackage{parskip}
  }{% else
    \setlength{\parindent}{0pt}
    \setlength{\parskip}{6pt plus 2pt minus 1pt}}
}{% if KOMA class
  \KOMAoptions{parskip=half}}
\makeatother
\usepackage{xcolor}
\usepackage[margin=1in]{geometry}
\usepackage{graphicx}
\makeatletter
\newsavebox\pandoc@box
\newcommand*\pandocbounded[1]{% scales image to fit in text height/width
  \sbox\pandoc@box{#1}%
  \Gscale@div\@tempa{\textheight}{\dimexpr\ht\pandoc@box+\dp\pandoc@box\relax}%
  \Gscale@div\@tempb{\linewidth}{\wd\pandoc@box}%
  \ifdim\@tempb\p@<\@tempa\p@\let\@tempa\@tempb\fi% select the smaller of both
  \ifdim\@tempa\p@<\p@\scalebox{\@tempa}{\usebox\pandoc@box}%
  \else\usebox{\pandoc@box}%
  \fi%
}
% Set default figure placement to htbp
\def\fps@figure{htbp}
\makeatother
\setlength{\emergencystretch}{3em} % prevent overfull lines
\providecommand{\tightlist}{%
  \setlength{\itemsep}{0pt}\setlength{\parskip}{0pt}}
\setcounter{secnumdepth}{-\maxdimen} % remove section numbering
\usepackage{bookmark}
\IfFileExists{xurl.sty}{\usepackage{xurl}}{} % add URL line breaks if available
\urlstyle{same}
\hypersetup{
  hidelinks,
  pdfcreator={LaTeX via pandoc}}

\author{}
\date{\vspace{-2.5em}}

\begin{document}

title: ``A codebook for collecting data from archaeological site forms
in Louisiana (v1.3)'' author: ``Jonathan Paige\footnote{New Mexico
  Consortium,
  \href{mailto:jonathan.n.paige@Gmail.com}{\nolinkurl{jonathan.n.paige@Gmail.com}}}''
output: html\_document: df\_print: paged pdf\_document: keep\_tex: yes
number\_sections: yes toc: true toc\_depth: 3 word\_document: default
header-includes: -

\usepackage{float}

\begin{itemize}
\item
  \floatplacement{figure}{hp}
\item
  \usepackage{booktabs}
\item
  \usepackage[font=small,labelfont=bf]{caption}
\item
  \usepackage{setspace}\doublespacing
\end{itemize}

\section{Version details:}\label{version-details}

This report was generated on 2026-10-05 08:00:04.929612

The current Git commit details are:

\begin{verbatim}
## Local:    main C:/Users/jpaige/Desktop/research_repositories/Site_form_codebook
## Remote:   main @ origin (https://github.com/jnpaige/Site_form_codebook.git)
## Head:     [79b2c9a] 2026-08-13: updating the frontmatter to support integration with codebook tools
\end{verbatim}

\section{This codebook}\label{this-codebook}

This codebook was developed as part of a broader project investigating
the reliability of large language models (LLMs) in performing text
coding tasks in archaeology. It is targeted specifically to site forms
on file with the Louisiana Division of Archaeology (Louisiana SHPO).
Traits include presence/absence judgements about cultural components and
artifact classes, as well as categorical and numeric extractions. These
judgements must be based on a reading of the narrative and structured
fields across site forms. This codebook provides guidance on making
those judgements.

\section{Source Structure}\label{source-structure}

Louisiana SHPO site forms contain a mix of structured fields
(checkboxes, dropdown selections, coded entries) and free-text narrative
sections. All are considered valid sources of evidence unless otherwise
specified in individual trait definitions below. Pre-printed field
labels on the form (e.g., column headers in artifact tallies) should not
be treated as evidence that a trait is present; only filled-in values,
checked boxes, or explicit narrative statements count as evidence.
Furthermore, the site forms are appended to one another when new work is
performed. In this context, we split the form into individual
investigations, and the relevant pages reporting that investigation, and
only provide those to the coder.The findings of an earlier investigation
may find different things from a later investigation. So, each
investigation should be coded independently. Prior to this coding
process, we broke the pdfs up into just pages that are from a particular
investigation, and which were likeliest to have relevant information
about the traits being coded.

\section{Definitions}\label{definitions}

\textbf{Pre-Contact period} Any occupation, use, or deposition of
materials at a site that pre-dates sustained European contact in the
relevant region of Louisiana. For most of Louisiana this is broadly
taken as pre-dating the late 17th century (ca. pre-1700 CE), though
recorders may use different cut-off dates. Both explicit statements and
inferred dates (e.g., radiocarbon dates, diagnostic artifact types) are
acceptable evidence.

\textbf{Post-Contact period} Any occupation, use, or deposition of
materials at a site that post-dates sustained European contact (broadly
ca. post-1700 CE in Louisiana). Includes colonial French, Spanish,
American, and later periods. Historic-era Euro-American, Creole, and
African-descended communities may all fall within this period.

\textbf{Caddo period} Occupation or material culture attributed to
Caddo-speaking peoples or their archaeological antecedents in Louisiana,
broadly spanning ca. 800--1700 CE (Late Woodland through Protohistoric),
though boundaries are debated and vary by subregion. Includes both the
Formative Caddo and Caddo tradition periods. Dating may be based on
diagnostic ceramics, architecture, or radiocarbon dates.

\textbf{Archaic period} Occupation or material culture attributed
broadly to the Archaic, spanning roughly 10,000--3,000 BP in Louisiana,
though exact boundaries vary by subregion and author. Includes Early,
Middle, and Late Archaic subperiods. Dating may be based on diagnostic
projectile point types, radiocarbon dates, or stratigraphic position.

\textbf{Chipped stone} Any artifact produced through knapping
(percussion or pressure flaking), including debitage (flakes, cores,
shatter), tools (projectile points, scrapers, bifaces, etc.), and
partially reduced nodules. Groundstone tools that also bear flaking
scars do not automatically qualify; the primary reduction method should
be knapping.

\textbf{Projectile point} A bifacially or unifacially retouched tool
interpreted as a tip or hafted element of a projectile weapon (spear,
atlatl dart, or arrow). Formal typological assignments (e.g., ``Gary
point'', ``Perdiz point'') are sufficient evidence. Fragments described
as probable projectile points are acceptable. Drills and perforators
should not be coded as projectile points unless the recorder explicitly
identifies them as such.

\textbf{Fire-cracked rock (FCR)} Thermally fractured rock resulting from
exposure to heat, typically associated with cooking features, earth
ovens, or hearths. May be described as burned rock, fire-cracked rock,
FCR, heat-fractured rock, or similar terms. Rock that is merely
discolored or scorched but not fractured does not qualify.

\textbf{NRHP eligibility} Status of a site with regard to National
Register of Historic Places inclusion. Eligible sites are eligible for
inclusion; listed sites are listed; most sites may have undetermined or
potential eligibility, barring more information typically gained through
historical records or more intensive fieldwork focused on testing the
site for research potential.

\section{Traits}\label{traits}

\subsection{Pre-Contact component
(Pre\_Conta)}\label{pre-contact-component-pre_conta}

\textbf{Data type:} binary

\textbf{Short Description:} Evidence of a pre-contact occupation or use
of the site.

\textbf{Definition:} The site form records evidence either in structured
fields or narrative text indicating that some portion of the site's
archaeological record dates to the period before sustained European
contact (broadly pre-1700 CE in Louisiana). This may include diagnostic
artifacts, radiocarbon dates, or interpretive statements by the
recorder.

\textbf{Inclusion criteria:} Evidence in text OR indicated in structured
fields (e.g., period checkboxes, cultural affiliation fields, artifact
type tallies referencing prehistoric types).

\textbf{Exclusion criteria:} Not described in text AND not indicated in
structured fields. A site described as exclusively historic-era should
be coded as absent.

\textbf{Typical exemplars:} Recorder checks ``prehistoric'' in the
cultural period field. Narrative states cases like ``the site contains a
prehistoric component'' or a cultural affiliation field lists ``unknown
prehistoric.'' Sites that are lithic scatters, or map onto any cultural
phase before European contact.

\textbf{Atypical exemplars:} Historic sites with presence of prehistoric
artifacts. A site is described as ``multicomponent'' and the prehistoric
component is mentioned only in passing in the survey narrative.

\textbf{Close but no:} A historic-era site form that mentions ``no
prehistoric materials observed.'' Presence of burned rock or undated
features alone, with no other indicator of pre-contact occupation. A
site form for an extant historic structure that references ``prehistoric
regional context'' in a background section without implying materials
were recovered.

\textbf{Response format:} Binary. Return \texttt{true} if the trait is
present, \texttt{false} if absent.

\subsection{Post-Contact component
(Post\_Conta)}\label{post-contact-component-post_conta}

\textbf{Data type:} binary

\textbf{Short Description:} Evidence of a post-contact occupation or use
of the site.

\textbf{Definition:} The site form records evidence indicating that some
portion of the site's archaeological record dates to the period after
sustained European contact (broadly post-1700 CE in Louisiana). This may
include historic-era artifact types, documentary references,
architectural remains, or explicit interpretive statements.

\textbf{Inclusion criteria:} Described in text OR indicated in
structured fields (e.g., period checkboxes listing ``historic,''
``colonial,'' ``19th century,'' etc.; artifact tallies listing historic
ceramics, glass, nails, or other temporally diagnostic historic
materials).

\textbf{Exclusion criteria:} Not described in text AND not indicated in
structured fields. A site described as exclusively prehistoric should be
coded as absent.

\textbf{Typical exemplars:} Recorder checks ``historic'' in the period
field. Narrative states ``the site contains historic-era debris
including European-made ceramics and cut nails.'' Cultural affiliation
listed as ``Anglo-American'' or ``French Colonial.'' A standing
structure on the site dates to the 19th century.

\textbf{Atypical exemplars:} No explicit period designation, but the
artifact inventory lists whiteware sherds, window glass, and wire nails
with no mention of prehistoric materials. A site described as
``multicomponent'' where the historic component is noted only briefly.

\textbf{Close but no:} A prehistoric site whose background section
discusses the broader historic settlement of the region. Presence of
modern intrusions (e.g., fence posts, drainage pipes) noted but not
interpreted as a historic-era occupation. A site where the recorder
mentions a nearby historic structure but does not indicate association
with the recorded site.

\textbf{Response format:} Binary. Return \texttt{true} if the trait is
present, \texttt{false} if absent.

\subsection{Caddo period component
(Caddo\_Unkn)}\label{caddo-period-component-caddo_unkn}

\textbf{Data type:} binary

\textbf{Short Description:} Evidence of a Caddo-period occupation,
regardless of specific subperiod assignment.

\textbf{Definition:} The site form records evidence attributing some
portion of the site's archaeological record to Caddo-speaking peoples or
the Caddo archaeological tradition, broadly spanning ca. 800--1700 CE.
The code name ``Caddo\_Unkn'' reflects that the specific Caddo subperiod
(e.g., Early, Middle, Late Caddo; Formative Caddo; Protohistoric) may be
unknown or unspecified. This trait should be coded as present whenever
Caddo affiliation is indicated, regardless of whether a finer subperiod
is also assigned.

\textbf{Inclusion criteria:} Described in text OR indicated in
structured fields (e.g., cultural affiliation field lists ``Caddo,''
``Caddoan,'' or a named Caddo-tradition phase; artifact field lists
Caddo ceramic types).

\textbf{Exclusion criteria:} Not described in text AND not indicated in
structured fields. Archaic, Woodland, or other non-Caddo prehistoric
components should not be coded as present here.

\textbf{Typical exemplars:} Cultural affiliation field states ``Caddo.''
Narrative describes ``Caddo ceramic sherds including engraved fine
ware.'' A named phase such as ``Haley,'' ``Belcher,'' or ``Kinsloe'' is
listed. The recorder assigns the site to the ``Late Prehistoric'' period
and describes Caddo-tradition ceramic types.

\textbf{Atypical exemplars:} The recorder uses only the term ``Late
Prehistoric'' without naming Caddo explicitly, but the artifact
inventory includes ceramic types (e.g., Friendship Engraved, Hickory
Fine Engraved) diagnostic of the Caddo tradition. A site recorded near
known Caddo territory where the artifact assemblage is consistent with
Caddo occupation but not explicitly labeled as such by the recorder.

\textbf{Close but no:} Sites described as Woodland-period only, even if
geographically within the Caddo region. Sites with a general
``prehistoric'' designation and no specific Caddo indicators in either
the structured fields or the narrative.

\textbf{Response format:} Binary. Return \texttt{true} if the trait is
present, \texttt{false} if absent.

\subsection{Archaic period component
(Arch\_Unkn)}\label{archaic-period-component-arch_unkn}

\textbf{Data type:} binary

\textbf{Short Description:} Evidence of an Archaic-period occupation,
regardless of specific subperiod (Early, Middle, or Late Archaic).

\textbf{Definition:} The site form records evidence attributing some
portion of the site's archaeological record to the Archaic period,
broadly spanning ca. 10,000--3,000 BP. The code name ``Arch\_Unkn''
reflects that the specific Archaic subperiod may be unknown. This trait
is coded as present whenever an Archaic occupation is indicated,
regardless of whether a finer temporal assignment is also made.

\textbf{Inclusion criteria:} Described in text OR indicated in
structured fields (e.g., period field lists ``Archaic,'' ``Early
Archaic,'' ``Middle Archaic,'' ``Late Archaic''; artifact field lists
Archaic diagnostic projectile point types).

\textbf{Exclusion criteria:} Not described in text AND not indicated in
structured fields. Paleoindian, Woodland, Caddo, or other non-Archaic
components should not be coded as present here.

\textbf{Typical exemplars:} Period field states ``Archaic.'' Narrative
describes ``a Late Archaic midden.'' Diagnostic projectile point types
associated with the Archaic (e.g., Gary, Motley, Evans, Pontchartrain)
are listed in the artifact inventory. Recorder notes ``undifferentiated
Archaic lithic scatter.''

\textbf{Atypical exemplars:} No explicit period designation, but the
artifact inventory lists only Gary or Evans points (well-established
Archaic diagnostics in Louisiana) with no other temporal indicators. A
site described as a shell midden that the recorder assigns to
``prehistoric'' but whose radiocarbon dates fall within the Archaic
range.

\textbf{Close but no:} A site described as ``prehistoric'' with only
generic lithic debitage and no diagnostic materials. Sites assigned to
the Woodland period only. Sites where the recorder mentions ``possible
Archaic materials'' without further support. Paleoindian components
only.

\textbf{Response format:} Binary. Return \texttt{true} if the trait is
present, \texttt{false} if absent.

\subsection{Chipped stone
(Chipped\_Stone)}\label{chipped-stone-chipped_stone}

\textbf{Data type:} binary

\textbf{Short Description:} Presence of chipped stone artifacts,
including tools and/or debitage.

\textbf{Definition:} The site form records the presence of any artifact
produced primarily through knapping (flaking), including but not limited
to: debitage (flakes, cores, angular shatter), bifaces, projectile
points, scrapers, drills, perforators, and other flaked stone tools.
Groundstone tools are excluded unless they also bear unambiguous
knapping scars interpreted as intentional reduction.

\textbf{Inclusion criteria:} Described in text OR indicated in
structured fields (e.g., artifact tally includes debitage, flakes,
cores, bifaces, or other knapped tool categories; narrative mentions
lithic scatter or chipped stone).

\textbf{Exclusion criteria:} Not described in text AND not indicated in
structured fields. Sites with only groundstone, ceramics, or other
non-knapped artifact classes should be coded as absent.

\textbf{Typical exemplars:} Artifact tally includes a count of ``lithic
debitage'' or ``flakes.'' Narrative states ``chipped stone tools were
recovered.'' Projectile points, scrapers, or bifaces are listed in the
artifact inventory. The site is described as a ``lithic scatter.''

\textbf{Atypical exemplars:} No explicit mention of chipped stone, but a
projectile point is listed in the artifact inventory (projectile points
are chipped stone by definition). Narrative mentions ``prehistoric
artifacts'' and the only illustrated or described items are flaked
tools.

\textbf{Close but no:} Only groundstone (e.g., grinding stones, manos,
metates) present, with no mention of knapped materials. Only ceramics
and historic-era artifacts listed. The word ``stone'' appears only in
the context of construction materials or geology with no artifact
referent.

\textbf{Response format:} Binary. Return \texttt{true} if the trait is
present, \texttt{false} if absent.

\subsection{Projectile point
(Proj\_Point)}\label{projectile-point-proj_point}

\textbf{Data type:} binary

\textbf{Short Description:} Presence of one or more projectile points.

\textbf{Definition:} The site form records the presence of at least one
artifact identified as a projectile point --- a bifacially or
unifacially retouched tool interpreted as the tip or hafted element of a
spear, atlatl dart, or arrow. Formal typological designations (e.g.,
``Gary point,'' ``Alba point,'' ``Perdiz point'') are sufficient.
Fragments described as probable projectile points are acceptable. Tools
identified primarily as drills or perforators are excluded unless the
recorder explicitly also identifies them as projectile points.

\textbf{Inclusion criteria:} Described in text OR indicated in
structured fields (e.g., artifact tally includes ``projectile point,''
``point,'' or a named type; narrative describes points or arrowheads).

\textbf{Exclusion criteria:} Not described in text AND not indicated in
structured fields. General lithic debitage or undifferentiated flakes
without mention of projectile points should be coded as absent.

\textbf{Typical exemplars:} Artifact tally lists ``3 projectile
points.'' Narrative states ``a Gary point was recovered from the
surface.'' The recorder lists ``Alba points'' or ``arrowheads'' in the
artifact inventory. A point fragment is described as such.

\textbf{Atypical exemplars:} No explicit category labeled ``projectile
point,'' but the artifact inventory lists a named type that is
unambiguously a projectile point form (e.g., ``2 Perdiz,'' ``1
Scallorn''). Narrative describes ``a bifacial point tip fragment.''

\textbf{Close but no:} Only debitage or cores listed with no mention of
finished point forms. A ``possible projectile point'' noted by the
recorder without further description. Drills or perforators listed but
not described as projectile points. The word ``point'' appears only in a
geographic context (e.g., ``Poverty Point culture'' or a place name).
Descriptions of bifaces, biface preforms, or blanks without further
specification of whether they represent projectile points.

\textbf{Response format:} Binary. Return \texttt{true} if the trait is
present, \texttt{false} if absent.

\subsection{Fire-cracked rock (FCR)}\label{fire-cracked-rock-fcr}

\textbf{Data type:} binary

\textbf{Short Description:} Presence of fire-cracked rock (thermally
fractured rock).

\textbf{Definition:} The site form records the presence of rock that has
been fractured as a result of exposure to high heat, typically
associated with cooking features, earth ovens, hearths, or burned rock
middens. May be referred to as ``fire-cracked rock,'' ``FCR,'' ``burned
rock,'' ``heat-fractured rock,'' ``fire-altered rock,'' or similar
terms. Rock that is merely discolored or stained without fracturing does
not qualify.

\textbf{Inclusion criteria:} Described in text OR indicated in
structured fields (e.g., artifact or feature tally includes FCR, burned
rock, or a feature type associated with FCR such as ``burned rock
midden'' or ``earth oven'').

\textbf{Exclusion criteria:} Not described in text AND not indicated in
structured fields. Thermally discolored but non-fractured rock does not
qualify. Burned structural materials (e.g., burned daub, burned brick)
are excluded.

\textbf{Typical exemplars:} Artifact tally includes a count or weight of
``FCR'' or ``fire-cracked rock.'' Narrative states ``burned rock was
recovered from a hearth feature.'' A feature is described as a ``burned
rock midden.'' The recorder notes ``considerable fire-cracked
limestone.''

\textbf{Atypical exemplars:} The site is described as containing hearth
features and the artifact inventory includes ``burned material'' if it
is clear from context that this refers to rock rather than organic
material or ceramics, code as present. A burned rock cluster is
described in the feature section but not explicitly tallied in the
artifact section.

\textbf{Close but no:} Only burned bone, burned ceramics, or burned
organic material mentioned, with no reference to burned or fractured
rock. Scorched sediment or ash described without associated fractured
rock. A hearth feature is listed but no mention is made of associated
rock. ``Burned'' applied only to structural daub or construction
materials.

\textbf{Response format:} Binary. Return \texttt{true} if the trait is
present, \texttt{false} if absent.

\subsection{NRHP Eligibility}\label{nrhp-eligibility}

\textbf{Data type:} categorical

\textbf{Short Description:} Status of site as a potential candidate for
inclusion in the National Register of Historic places or as a protected
resource.

\textbf{Definition:} Eligibility statement outlining the current status
of the site, or the recommendation for the site.

\textbf{Inclusion criteria:} Code if National Register eligibilty, site
eligibility, research potential of the site as it relates to NRHP
eligibility, or presence of graves or other protected resources is
mentioned.

\textbf{Exclusion criteria:} No discussion about national register
eligibility, site eligibility, research potential, or protected status
of the site.

\textbf{Categories:}

Eligible: The site is judged as eligible for inclusion.

Listed: The site had been listed in the register.

Potential/undetermined eligibility: The site is marked as having
potential for inclusion in the register, or there is a note that more
research is needed to determine eligibility status. This typically is
paired with a recommendation for further study.

Ineligible: The site is marked as not being eligible for inclusion

Protected: The site is noted as not being eligible for inclusion, but is
otherwise a protected resource, like a cemetery or homestead with
graves.

Unknown: The particular site form, or record of study in a particular
field campaign, has no information about recommendations for NRHP
eligibility.

\textbf{Response format:} Categorical/character entry containing one of
the following: ``Eligible'', ``Listed'', ``Potential/undetermined
eligibility'', ``Ineligible'', ``Protected'', ``Unknown''.

\subsection{Site depth}\label{site-depth}

\textbf{Data type:} numeric/character

\textbf{Short Description:} The depth of the deepest sediments at a site
that contain archaeological material.

\textbf{Definition:} The field identified depth of archaeological
material at a site, typically identified through either shovel testing
or excavation of 1x1m units.

\textbf{Inclusion criteria:} Descriptions of excavation results are
present in structured fields or in narrative text, and those
descriptions include discussion of buried sediments that contain
archaeological artifacts or features.

\textbf{Exclusion criteria:} Site is an unexcavated site that has
largely architectural or surface features.

\textbf{Response format:} Return a plain number in cm below surface
(e.g., \texttt{40} not \texttt{"40\ cmbs"} or \texttt{"0.4\ m"}).
Convert to cm if reported in other units. If site was excavated, but
there is not enough information to determine the depth of archaeological
deposits, return ``not reported''. If no excavation occurred, return
\texttt{"unknown"}.

\subsection{Site area}\label{site-area}

\textbf{Data type:} numeric/character

\textbf{Short Description:} The reported area of the archaeological
site.

\textbf{Definition:} The field identified area of the site, either in
terms of two dimensions (e.g.~the length of the site on the East West
axis, and the length of the site on the North South axis), or as an
areal measurement.

\textbf{Inclusion criteria:} The area of the site is described either in
narrative, or standardized fields.

\textbf{Exclusion criteria:} The site is a landscape scale feature,
network of features (like historic roads or trails), or otherwise linear
features that have no clearly reported area.

\textbf{Response format:} Return a number in square meters but do not
include the unit (e.g.~return: 450. Do not return: 450m\^{}2. Convert to
square meters if area is reported in other units, like square
kilometers, acres, or hectares. If only dimensions are reported,
calculate area from those dimensions assuming a rectlinear shape.

\subsection{Methodology}\label{methodology}

\textbf{Data type:} list

\textbf{Short Description:} Methodologies employed during the site visit
reported in the form.

\textbf{Definition:} Any one or more methodologies performed during the
site visit.

\textbf{Inclusion criteria:} Some kind of site visit, identification, or
assessment is described.

\textbf{Exclusion criteria:} By definition, any of the site forms should
have some form of methodology reported, at minimum the site was visited.

\textbf{Categories:}

Surface survey: Site evaluate through surface survey

Shovel testing: Site evaluated through excavating one or more shovel
test units

Test unit: Site evaluated through hand excavated test units, often in
1x1 meter cells

Backhoe trenching: Site evaluated through excavation of trenches by a
backhoe or some other form of construction machinery

Architectural survey: Architectural assessment performed on features at
the site

Historical records/interviews: Information about the site gathered
through historical records or local informants

Geological: Formal geoarchaeological assessment

\textbf{Response format:} JSON array containing zero or more of:
``Surface survey'', ``Shovel testing'', ``Test unit'', ``Backhoe
trenching'', ``Architectural survey'', ``Historical
records/interviews'', ``Geological''

\subsection{Tramline or rail
infrastructure}\label{tramline-or-rail-infrastructure}

\textbf{Data type:} binary

\textbf{Short Description:} Remains of rail-based transport, including
logging tram grades and railroad earthworks.

\textbf{Definition:} The form identifies remains of a rail or tram
corridor --- a tram grade or railroad berm, flanking borrow ditches,
cuts through rising ground.

\textbf{Inclusion criteria:} Described in text OR indicated in
structured fields. Code present when the recorder names a tram,
tramline, railroad, rail bed, rail grade, or spur, or describes a linear
feature and argues that it served rail or tram transport.

\textbf{Exclusion criteria:} Not described in text AND not indicated in
structured fields. A linear raised feature, berm, or parallel ditch
system recorded without the recorder attributing it to rail or tram use.
Linear features the form attributes to another function --- a levee, a
drainage or irrigation ditch, a road, a logging skid trail, a pipeline
or utility corridor, a firebreak.

\textbf{Typical exemplars:} ``This site includes a tramline feature
running NE/SW for 30 meters''

\textbf{Atypical exemplars:} ``The site is disturbed by a historic
tramline which cut the site at the NE corner''

\textbf{Close but no:} Berms and grades not described as relating to a
tramline.

\textbf{Response format:} Binary. Return \texttt{true} if the trait is
present, \texttt{false} if absent.

\subsection{Other logging
infrastructure}\label{other-logging-infrastructure}

\textbf{Data type:} binary

\textbf{Short Description:} archaeological evidence for logging
independent of tramlins: mills, camps, equipment, artifacts.

\textbf{Definition:} The form describes archaeological remains of
commercial timber extraction or processing that are not tramline or rail
features. This covers remains of sawmill complexes and their footings,
boilers, mill ponds, splash dams, and sawdust or slab deposits; logging
camps and workers' quarters; skidder sets, loading decks, and log
landings; and the discarded hardware of those operations, such as cable,
steam-donkey parts, and saw fragments.

\textbf{Inclusion criteria:} Described in text OR indicated in
structured fields. Code present when authors argue that the findings
relate to the logging industry --- individual artifacts like logging
cable, landscape features like a levelled mill or camp pad, or an
encampment the recorder reads as a logging camp.

\textbf{Exclusion criteria:} Not described in text AND not indicated in
structured fields. Tramline and rail features, which are coded under
Tramline or rail infrastructure. Historic refuse with no stated or
arguable connection to timber work --- a can scatter, however dense,
that the recorder does not interpret as a logging camp, mill, or other
logging feature.

\textbf{Typical exemplars:} The recorder describes a levelled bench with
a tight concentration of early-twentieth-century cans and argues it
represents a logging camp, although no foundations survive.

\textbf{Atypical exemplars:} Concrete sawmill footings with an adjacent
mill pond and slab pile. A logging camp with scattered domestic refuse,
a dense can dump, and structural piers. Steam-donkey or skidder cable
and parts.

\textbf{Close but no:} A tram berm with no other logging evidence
described. Description of impacts to sites due to modern logging
activities.

\textbf{Response format:} Binary. Return \texttt{true} if the trait is
present, \texttt{false} if absent.

\subsection{Turpentine cups}\label{turpentine-cups}

\textbf{Data type:} binary

\textbf{Short Description:} Clay or metal cups used to collect resin
from catfaced pines.

\textbf{Definition:} The form describes turpentine cups --- the
Herty-system terracotta cups introduced in the early twentieth century,
or the later metal cups used to catch resin running from a catface on a
worked pine. Tin gutters or aprons nailed to the tree to direct resin
into the cup fall under this trait.

\textbf{Inclusion criteria:} Described in text OR indicated in
structured fields. Code present when the recorder records turpentine
cups, Herty cups, or resin cups by any of those names, or the gutters
and aprons associated with them. Fragments the recorder identifies as
cup sherds are sufficient.

\textbf{Exclusion criteria:} Not described in text AND not indicated in
structured fields. Undifferentiated historic ceramic or metal containers
the recorder does not identify as turpentine-related. Catfaces or boxes
with no cup described. Those are coded under ``Other turpentine
evidence''.

\textbf{Typical exemplars:} ``Several Herty cups were observed at the
bases of catfaced pines.'' Terracotta turpentine cup fragments in a
surface scatter. Metal resin cups and tin gutters collected.

\textbf{Atypical exemplars:} A scatter of coarse terracotta sherds the
recorder identifies as cup fragments.

\textbf{Close but no:} Flower pot sherds. Undiagnostic coarse
earthenware..

\textbf{Response format:} Binary. Return \texttt{true} if the trait is
present, \texttt{false} if absent.

\subsection{Tar kiln}\label{tar-kiln}

\textbf{Data type:} binary

\textbf{Short Description:} Potential earth kiln for rendering tar and
pitch from fat lightwood.

\textbf{Definition:} The form describes a tar kiln: a broad, shallow,
roughly circular depression, commonly 5 to 15 m across, where resinous
lightwood was slow-burned under an earth cover to render tar and pitch.
Charcoal, ash, or burned sand is frequently noted, and sometimes a
prepared clay floor or a central collection channel running to a
catchment at the edge. Related terms include pitch kiln, tar pit, and
tar kiln scar.

\textbf{Inclusion criteria:} Described in text OR indicated in
structured fields. Code present when the recorder names a tar kiln,
pitch kiln, or tar pit.

\textbf{Exclusion criteria:} Not described in text AND not indicated in
structured fields. Pits the recorder attributes to gum turpentine
production. Depressions the recorder attributes to another cause --- a
cellar, privy, well, borrow pit, treefall, or stock pond --- or records
without interpreting.

\textbf{Typical exemplars:} ``A tar kiln approximately 8 m in diameter
with a charcoal-stained floor is present on the ridge nose.'' The
recorder maps and names a pitch kiln.

\textbf{Atypical exemplars:} A broad shallow circular depression with
charcoal flecking that the recorder interprets as probable tar
production without naming a kiln type.

\textbf{Close but no:} A rectangular depression interpreted as a
structural cellar. A deep narrow pit. A modern burn pile. A turpentine
still firebox. A broad shallow depression the recorder records but does
not interpret. Such depressions are more appropriately coded as
``Unattributed broad pit''.

\textbf{Response format:} Binary. Return \texttt{true} if the trait is
present, \texttt{false} if absent.

\subsection{Unattributed broad pit}\label{unattributed-broad-pit}

\textbf{Data type:} binary

\textbf{Short Description:} A broad shallow depression of tar-kiln form
that the recorder describes but does not interpret.

\textbf{Definition:} The form describes a pit or depression whose
morphology is consistent with a tar kiln, broad, shallow, roughly
circular, on the order of several metres across, but assigns it no
function, or records its function as unknown. A broad shallow depression
may be a tar kiln, a borrow pit, a collapsed cellar, a treefall, or a
natural hollow.

\textbf{Inclusion criteria:} Described in text OR indicated in
structured fields. Code present when the form describes a broad,
shallow, roughly circular depression and either assigns it no function
or records the function as unknown, uncertain, or undetermined.
Measurements help but are not required where the description conveys the
form.

\textbf{Exclusion criteria:} Not described in text AND not indicated in
structured fields. Any depression the recorder identifies as a tar kiln,
pitch kiln, or tar pit --- those belong to Tar kiln. Any depression the
recorder confidently attributes to another function: a structural
cellar, privy, well, borrow pit, stock pond, treefall, or crater. Narrow
or deep pits, rectangular depressions, and linear ditches, whose form is
not that of a tar kiln.

\textbf{Typical exemplars:} ``A shallow circular depression
approximately 7 m in diameter is present on the ridge; function
unknown.'' ``Feature 2 is a broad basin-shaped depression of
undetermined origin.''

\textbf{Atypical exemplars:} A depression recorded only on the sketch
map with a dimension and no label, where the narrative does not discuss
it. A recorder noting ``several shallow depressions of uncertain
origin'' without describing them individually.

\textbf{Close but no:} A depression the recorder calls a probable borrow
pit. A cellar hole. A treefall the recorder identifies as such. A deep
narrow shaft. A depression whose dimensions are not given and whose
description does not convey a broad, shallow form.

\textbf{Response format:} Binary. Return \texttt{true} if the trait is
present, \texttt{false} if absent.

\subsection{Other turpentine evidence}\label{other-turpentine-evidence}

\textbf{Data type:} binary

\textbf{Short Description:} Naval-stores evidence that is neither a cup
nor a pit.

\textbf{Definition:} The form describes evidence of turpentining or
naval-stores production other than turpentine cups and tar kilns. This
includes catface scars and chipped boxes on living pines or stumps;
still and distillery remains; cooperage and barrel hoops; the hand tools
of the trade, such as hacks, pullers, scrape irons, and dip buckets; and
camps or quarters the recorder attributes to turpentine operations.

\textbf{Inclusion criteria:} Described in text OR indicated in
structured fields. Code present when the recorder records naval-stores
evidence falling outside the cup and pit traits, or argues that a
feature or artifact relates to turpentine production.

\textbf{Exclusion criteria:} Not described in text AND not indicated in
structured fields. Evidence already captured by Turpentine cups or Tar
kiln. Historic material the recorder does not connect to naval stores.

\textbf{Typical exemplars:} ``Numerous catfaced longleaf pines are
present across the site.'' Remains of a turpentine still. Barrel hoops
and a scrape iron the recorder attributes to turpentining. A quarters
site the recorder assigns to a turpentine camp.

\textbf{Atypical exemplars:} A stand of stumps bearing paired diagonal
scars that the recorder interprets as boxing, without using the word
turpentine.

\textbf{Close but no:} Fire-scarred pines the recorder attributes to
burning. Undiagnostic barrel hardware with no stated naval-stores
association. A tar kiln, which is coded under Tar kiln.

\textbf{Response format:} Binary. Return \texttt{true} if the trait is
present, \texttt{false} if absent.

\subsection{Homestead or farmstead}\label{homestead-or-farmstead}

\textbf{Data type:} binary

\textbf{Short Description:} Remains of a rural domestic occupation, with
or without associated agricultural features.

\textbf{Definition:} The form describes the remains of a dwelling and
its associated yard and outbuildings. The form may also simply describe
scatters of artifacts, and the author may intuit that it is a
farmstead/homestead based on positioning on the landscape and other
context clues.

\textbf{Inclusion criteria:} Described in text OR indicated in
structured fields. Code present when the form describes domestic
structural remains, or a domestic refuse scatter in direct association
with a structure, well, or cistern. Also code present where the recorder
argues for a house site. Terms such as ``old home place'',
``homestead'', or ``farmstead'' are sufficient on their own.

\textbf{Exclusion criteria:} Not described in text AND not indicated in
structured fields. A refuse dump with no structural or water-feature
association and no argument by the recorder for a dwelling. Worker
housing within a logging camp, unless the form describes independent
domestic or agricultural occupation.

\textbf{Typical exemplars:} A chimney fall with surrounding whiteware,
bottle glass, and cut nails. Brick piers, a hand-dug well, and a privy
depression. ``The site is the location of the old Dowden home place.'' A
barn foundation with associated fencing.

\textbf{Atypical exemplars:} A levelled pad on a low rise with a bounded
concentration of whiteware, bottle glass, and nails, which the recorder
interprets as a house site despite no surviving foundation. A domestic
refuse scatter with an associated hand-dug well and no other structural
evidence.

\textbf{Close but no:} An isolated bottle dump in a borrow pit. A
logging camp described purely as temporary quarters. A standing modern
residence recorded as a land-use condition rather than as an
archaeological site.

\textbf{Response format:} Binary. Return \texttt{true} if the trait is
present, \texttt{false} if absent.

\end{document}
